\section{User Stories}
\label{chap:userStories}

Di seguito vengono presentate le User Story del progetto \textbf{WattGuard}, sviluppate a partire dai requisiti funzionali individuati nei capitoli precedenti.

\subsection{User Story 0 – Cambio lingua dell'interfaccia}
\label{us:cambioLingua}

\textbf{Selezione della lingua e localizzazione dell'interfaccia} \\
Come operatore tecnico o amministratore comunale, voglio poter selezionare la lingua dell'interfaccia tra italiano, inglese e tedesco, in modo da utilizzare il sistema nella mia lingua madre e facilitare la comprensione di tutti i contenuti.

\paragraph*{Criteri di accettazione:}
\begin{itemize}
    \item Il sistema supporta tre lingue: italiano (IT), inglese (EN) e tedesco (DE).
    \item È presente un selettore lingua (ad esempio, un menu a tendina o bandierine) in una posizione facilmente accessibile in tutte le pagine dell'applicazione.
    \item La selezione della lingua viene memorizzata e persiste tra le sessioni dell'utente.
    \item Tutte le stringhe visibili all'utente (etichette, messaggi, pulsanti, notifiche, ecc.) vengono tradotte nella lingua selezionata.
    \item I contenuti generati lato server (email, report, notifiche) rispettano la lingua selezionata dall'utente.
    \item Date, numeri e formati orari seguono le convenzioni della locale selezionata.
    \item La lingua di riferimento predefinita è l'inglese (EN).
\end{itemize}

\paragraph*{TASKS – User Story 0:}
\begin{enumerate}
    \item \textbf{Configurare il sistema di localizzazione}
    \begin{itemize}
        \item Integrare una libreria di internazionalizzazione nel frontend.
        \item Creare i cataloghi di traduzione per le tre lingue supportate (IT, EN, DE).
        \item Organizzare i file di traduzione per moduli/componenti per facilitarne la manutenzione.
    \end{itemize}

    \item \textbf{Implementare il selettore lingua nell'interfaccia}
    \begin{itemize}
        \item Aggiungere un controllo di selezione lingua (dropdown o bandierine) nell'header dell'applicazione.
        \item Assicurare che il selettore sia visibile in tutte le pagine protette e nella pagina di login.
        \item Gestire il cambio lingua con ricarica dinamica dell'interfaccia.
    \end{itemize}

    \item \textbf{Implementare la persistenza della preferenza lingua}
    \begin{itemize}
        \item Memorizzare la lingua selezionata nel localStorage o in un cookie.
        \item Recuperare la preferenza al successivo accesso dell'utente.
        \item Se l'utente non ha una preferenza salvata, utilizzare la lingua di default (EN) o rilevare quella del browser.
    \end{itemize}

    \item \textbf{Adattare il backend per la localizzazione}
    \begin{itemize}
        \item Configurare il backend per rilevare la lingua richiesta.
        \item Localizzare i contenuti generati dal server (email di invito, notifiche, report PDF/Excel).
        \item Formattare date, numeri e orari secondo la locale selezionata.
    \end{itemize}

    \item \textbf{Verificare e testare la localizzazione}
    \begin{itemize}
        \item Testare il cambio lingua in tutte le sezioni dell'applicazione.
        \item Verificare la corretta visualizzazione dei caratteri speciali (es. ä, ö, ü per il tedesco).
        \item Testare la persistenza della preferenza tra sessioni diverse.
        \item Verificare che email e report vengano generati nella lingua corretta.
        \item Testare il comportamento su dispositivi diversi (desktop, tablet).
    \end{itemize}
\end{enumerate}

\subsection{User Story 1 – Autenticazione e gestione sessione (Login e Logout)}
\label{us:login}

\textbf{Accesso flessibile e terminazione sicura della sessione} \\
Come operatore tecnico o amministratore comunale, voglio poter accedere al sistema tramite credenziali locali o autenticazione federata Google ed effettuare il logout al termine delle attività, in modo da operare secondo le mie autorizzazioni e garantire la sicurezza dei dati.

\paragraph*{Criteri di accettazione:}
\begin{itemize}
    \item Entrambi i ruoli (operatori e amministratori) possono scegliere se autenticarsi inserendo email e password oppure cliccando sul pulsante ``Accedi con Google''.
    \item A seguito di un'autenticazione riuscita, il sistema reindirizza l'utente alla dashboard generale comune.
    \item Se il login fallisce (per credenziali errate, account inesistente o errore del provider esterno), il sistema mostra un messaggio di errore chiaro e non bloccante.
    \item L'utente autenticato può terminare la sessione in qualsiasi momento tramite il comando di logout presente nel menù profilo; l'operazione rimuove il token di accesso dal dispositivo, ripulisce lo stato locale e reindirizza alla schermata di login.
    \item Il sistema gestisce la scadenza automatica della sessione (validità del token di accesso, 8 ore).
    \item Tutte le fasi di autenticazione e scambio dati avvengono su protocollo cifrato HTTPS.
\end{itemize}

\paragraph*{TASKS – User Story 1:}
\begin{enumerate}
    \item \textbf{Integrazione delle opzioni di login}
    \begin{itemize}
        \item Integrare nella pagina di login sia il form per le credenziali locali sia il pulsante SSO Google.
        \item Gestire il recupero password tramite apposito link contestuale.
    \end{itemize}

    \item \textbf{Configurazione dell'autenticazione federata}
    \begin{itemize}
        \item Configurare l'integrazione con Google.
        \item Mappare i token restituiti dal provider con le utenze censite a sistema.
    \end{itemize}

    \item \textbf{Gestione del routing e della sessione}
    \begin{itemize}
        \item Reindirizzare l'utente alla dashboard generale post-autenticazione.
        \item Implementare la procedura di logout con rimozione del token di accesso dal dispositivo (localStorage) e redirect al login.
    \end{itemize}

    \item \textbf{Sicurezza delle credenziali}
    \begin{itemize}
        \item Cifrare le password memorizzate localmente tramite hashing robusto (bcrypt o Argon2 con salt individuale).
        \item Proteggere le rotte applicative impedendo l'accesso ad utenti non autenticati o disconnessi.
    \end{itemize}

    \item \textbf{Verifica e test}
    \begin{itemize}
        \item Testare il ciclo completo di login (locale e Google), fallimento credenziali e corretto funzionamento del logout (impossibilità di accedere alle aree protette dopo la disconnessione).
    \end{itemize}
\end{enumerate}

\subsection{User Story 2 – Visualizzazione dashboard generale}
\label{us:dashboardGenerale}

\textbf{Visualizzazione panoramica consumi energetici} \\
Come operatore tecnico o amministratore comunale, voglio visualizzare una dashboard generale con grafici aggregati dei consumi energetici, in modo da avere una visione d'insieme immediata delle performance degli edifici pubblici.

\paragraph*{Criteri di accettazione:}
\begin{itemize}
    \item La dashboard mostra grafici aggregati dei consumi totali
    \item Vengono visualizzati gli indicatori principali: sensori attivi e totali, allarmi attivi e consumi correnti di elettricità e gas
    \item La dashboard si carica entro 10 secondi per il 95\% delle richieste
    \item I dati vengono aggiornati automaticamente
\end{itemize}

\paragraph*{TASKS – User Story 2:}
\begin{enumerate}
    \item \textbf{Progettare il layout della dashboard}
    \begin{itemize}
        \item Creare wireframe con disposizione ottimale di grafici e indicatori
        \item Assicurare design responsive per tablet con schermo minimo di 10 pollici
    \end{itemize}
    \item \textbf{Implementare i grafici di visualizzazione dati}
    \begin{itemize}
        \item Sviluppare grafici per consumi aggregati utilizzando librerie appropriate
        \item Creare indicatori visivi per temperature e performance
    \end{itemize}
    \item \textbf{Implementare il sistema di aggiornamento automatico}
    \begin{itemize}
        \item Configurare refresh automatico dei dati ogni 90 secondi (intervallo configurabile)
    \end{itemize}
    \item \textbf{Ottimizzare le query al database}
    \begin{itemize}
        \item Assicurare tempi di risposta inferiori a 10 secondi
        \item Implementare indici appropriati sulle collezioni
    \end{itemize}
    \item \textbf{Testare le performance}
    \begin{itemize}
        \item Verificare i tempi di caricamento con carichi simultanei di 20 utenti
        \item Testare la responsività su diversi dispositivi
    \end{itemize}
\end{enumerate}


\subsection{User Story 3 – Ricerca edifici con criteri multipli}
\label{us:ricercaEdifici}

\textbf{Ricerca edifici con criteri multipli} \\
Come operatore tecnico o amministratore, voglio cercare gli edifici pubblici monitorati utilizzando nome, indirizzo o tipologia di edificio, in modo da trovare rapidamente l'edificio di mio interesse.

\paragraph*{Criteri di accettazione:}
\begin{itemize}
    \item L'utente può inserire uno o più criteri di ricerca
    \item Il sistema restituisce risultati che corrispondono a tutti i criteri inseriti
    \item La ricerca si completa entro 5 secondi
    \item I risultati mostrano nome edificio, indirizzo, tipologia impianto e stato operativo
\end{itemize}

\paragraph*{TASKS – User Story 3:}
\begin{enumerate}
    \item \textbf{Creare l'interfaccia di ricerca}
    \begin{itemize}
        \item Implementare campi di input per nome, indirizzo e tipologia di edificio
        \item Aggiungere pulsante di ricerca e reset
    \end{itemize}
    \item \textbf{Implementare la logica di ricerca}
    \begin{itemize}
        \item Sviluppare query che combinano i criteri inseriti con operatore AND
    \end{itemize}
    \item \textbf{Visualizzare i risultati}
    \begin{itemize}
        \item Creare una tabella o lista con i dettagli degli edifici trovati
        \item Includere tutti i campi richiesti (nome, indirizzo, superficie, tipologia, stato, ultimo aggiornamento)
    \end{itemize}
    \item \textbf{Ottimizzare le performance}
    \begin{itemize}
        \item Creare indici appropriati sulle colonne di ricerca
        \item Assicurare tempi di risposta inferiori a 5 secondi
    \end{itemize}
    \item \textbf{Gestire i casi limite}
    \begin{itemize}
        \item Mostrare messaggio appropriato quando non ci sono risultati
        \item Validare gli input prima della ricerca
    \end{itemize}
\end{enumerate}


\subsection{User Story 4 – Visualizzazione dettaglio edificio}
\label{us:dettaglioEdificio}

\textbf{Visualizzazione completa dati edificio} \\
Come operatore tecnico o amministratore, voglio visualizzare tutti i dettagli di un edificio selezionato inclusi dati anagrafici, misurazioni in tempo reale e grafici storici, in modo da analizzare i consumi energetici in differenti intervalli di tempo avendo una visione completa delle sue condizioni.

\paragraph*{Criteri di accettazione:}
\begin{itemize}
    \item La schermata mostra informazioni anagrafiche complete (nome, indirizzo, superficie, anno costruzione)
    \item Vengono visualizzati i dati in tempo reale (temperatura interna, esterna, consumo istantaneo)
    \item Sono presenti grafici storici dei consumi e temperature
    \item Viene mostrato lo stato di funzionamento dei sensori
    \item L'utente può selezionare diversi intervalli temporali per i grafici storici scegliendo tra intervalli predefiniti (ultimi 7, 30, 90 o 365 giorni) oppure tramite la scelta di due date distinte
\end{itemize}

\paragraph*{TASKS – User Story 4:}
\begin{enumerate}
    \item \textbf{Creare il layout della pagina dettaglio}
    \begin{itemize}
        \item Organizzare le informazioni in sezioni logiche
        \item Assicurare una gerarchia visiva chiara
    \end{itemize}
    \item \textbf{Implementare la visualizzazione dati anagrafici}
    \begin{itemize}
        \item Recuperare e mostrare tutti i dati dell'edificio dal database
        \item Formattare le informazioni in modo leggibile
    \end{itemize}
    \item \textbf{Sviluppare la sezione dati in tempo reale}
    \begin{itemize}
        \item Implementare visualizzazione dei valori correnti con aggiornamento automatico
        \item Utilizzare indicatori visivi appropriati (gauge, termometri, ecc.)
    \end{itemize}
    \item \textbf{Creare i grafici storici}
    \begin{itemize}
        \item Implementare grafici interattivi per consumi e temperature
        \item Permettere la selezione di intervalli temporali diversi
    \end{itemize}
    \item \textbf{Aggiornare dinamicamente i grafici}
    \begin{itemize}
        \item Implementare il refresh dei grafici al cambio di selezione
        \item Assicurare transizioni fluide tra le visualizzazioni
    \end{itemize}
    \item \textbf{Mostrare lo stato dei sensori}
    \begin{itemize}
        \item Visualizzare lista sensori con indicatori di stato (attivo/offline)
        \item Mostrare data ultimo aggiornamento per ogni sensore
    \end{itemize}
\end{enumerate}


\subsection{User Story 5 – Esportazione dati}
\label{us:esportazioneDati}

\textbf{Esportazione dati in formato CSV} \\
Come amministratore comunale, voglio esportare i dati dei consumi energetici in formato CSV, in modo da poterli analizzare con strumenti esterni o includere in report personalizzati.

\paragraph*{Criteri di accettazione:}
\begin{itemize}
    \item L'utente può selezionare uno o più edifici per l'esportazione
    \item L'utente può specificare il periodo temporale di interesse
    \item Il file CSV contiene le rilevazioni selezionate (timestamp, edificio, sensore, tipologia, valore, unità di misura)
    \item Il download del file avviene immediatamente dopo la generazione
    \item I dati sono formattati correttamente per l'importazione in fogli di calcolo
    \item Sono disponibili inoltre esportazioni in formato JSON (rilevazioni e backup del database) e report aggregati in PDF/Excel
\end{itemize}

\paragraph*{TASKS – User Story 5:}
\begin{enumerate}
    \item \textbf{Creare l'interfaccia di esportazione}
    \begin{itemize}
        \item Implementare form per selezione edifici e periodo temporale
        \item Aggiungere opzione per formato CSV
    \end{itemize}
    \item \textbf{Implementare la logica di generazione CSV}
    \begin{itemize}
        \item Estrarre i dati dal database secondo i criteri selezionati
        \item Formattare i dati in formato CSV standard
    \end{itemize}
    \item \textbf{Includere tutte le informazioni necessarie}
    \begin{itemize}
        \item Assicurare che timestamp, edificio, sensore, tipologia, valore e unità di misura siano presenti
        \item Aggiungere intestazioni colonne appropriate
    \end{itemize}
    \item \textbf{Gestire il download del file}
    \begin{itemize}
        \item Implementare il trigger di download automatico
        \item Nominare il file in modo descrittivo (wattguard-consumption-{start}-{end}.csv)
    \end{itemize}
    \item \textbf{Testare la funzionalità}
    \begin{itemize}
        \item Verificare la corretta formattazione del CSV
        \item Testare l'importazione in Excel e altri strumenti
        \item Testare la qualità di stampa
    \end{itemize}
\end{enumerate}


\subsection{User Story 6 – Gestione utenti (creazione account)}
\label{us:creazioneAccount}

\textbf{Creazione account operatore/amministratore tramite invito} \\
Come amministratore comunale, voglio creare nuovi inviti per gli operatori tecnici e amministratori, specificando l'email e il ruolo, in modo da gestire centralmente gli accessi al sistema.

\paragraph*{Criteri di accettazione:}
\begin{itemize}
    \item L'amministratore può accedere a una sezione dedicata alla gestione utenti
    \item È possibile inserire i dati richiesti (email e ruolo del nuovo utente)
    \item Il sistema invia un'email automatica all'utente con il link di invito univoco
    \item Al primo accesso l'utente imposta la propria password oppure accetta l'invito tramite Google
    \item Il sistema valida l'unicità dell'email
    \item L'amministratore può revocare gli inviti pendenti
\end{itemize}

\paragraph*{TASKS – User Story 6:}
\begin{enumerate}
    \item \textbf{Creare l'interfaccia di gestione utenti}
    \begin{itemize}
        \item Implementare form per l'inserimento dell'email e del ruolo dell'invitato
        \item Implementare l'elenco degli inviti e la loro revoca
    \end{itemize}
    \item \textbf{Implementare la logica di creazione dell'invito}
    \begin{itemize}
        \item Validare i dati inseriti (formato email, campi obbligatori)
        \item Verificare l'unicità dell'indirizzo email
    \end{itemize}
    \item \textbf{Generare il token di invito}
    \begin{itemize}
        \item Generare un token di invito sicuro, memorizzato come hash
        \item Associare al token l'email e il ruolo dell'invitato, con scadenza (7 giorni)
    \end{itemize}
    \item \textbf{Implementare l'invio email automatico}
    \begin{itemize}
        \item Configurare il template email con il link di invito e le istruzioni
        \item Inviare email all'indirizzo specificato
    \end{itemize}
    \item \textbf{Tracciare l'operazione}
    \begin{itemize}
        \item Registrare l'amministratore che ha creato l'invito
    \end{itemize}
\end{enumerate}


\subsection{User Story 7 – Gestione utenti (modifica ruolo)}
\label{us:modificaPermessi}

\textbf{Modifica ruolo operatore/amministratore} \\
Come amministratore comunale, voglio modificare il ruolo di un operatore tecnico/amministratore esistente, in modo da adattare i suoi accessi alle nuove responsabilità o esigenze organizzative.

\paragraph*{Criteri di accettazione:}
\begin{itemize}
    \item L'amministratore può visualizzare la lista degli utenti esistenti
    \item È possibile selezionare un utente e modificarne il ruolo (operatore/amministratore)
    \item Le modifiche vengono salvate immediatamente
    \item L'amministratore non può modificare il proprio ruolo
\end{itemize}

\paragraph*{TASKS – User Story 7:}
\begin{enumerate}
    \item \textbf{Creare la visualizzazione lista operatori}
    \begin{itemize}
        \item Implementare tabella con tutti gli operatori e il loro ruolo attuale
        \item Aggiungere pulsante di modifica per ogni operatore
    \end{itemize}
    \item \textbf{Implementare l'interfaccia di modifica ruolo}
    \begin{itemize}
        \item Creare selettore del ruolo (operatore/amministratore)
        \item Pre-popolare con il ruolo attuale dell'operatore
    \end{itemize}
    \item \textbf{Implementare la logica di aggiornamento}
    \begin{itemize}
        \item Validare le modifiche prima del salvataggio
        \item Aggiornare il database con il nuovo ruolo
    \end{itemize}
\end{enumerate}


\subsection{User Story 8 – Configurazione soglie di allarme}
\label{us:soglieAllarme}

\textbf{Definizione soglie personalizzate} \\
Come operatore tecnico o amministratore comunale, voglio configurare le soglie di allarme dei sensori (valore minimo e massimo, ad esempio per temperatura e consumi), in modo da ricevere segnalazioni solo per situazioni realmente anomale.

\paragraph*{Criteri di accettazione:}
\begin{itemize}
    \item L'utente può definire per ciascun sensore la soglia minima e massima del valore misurato
    \item Il sistema verifica la coerenza dei valori (soglia minima inferiore alla massima)
    \item Le soglie possono essere rimosse dal sensore quando non più necessarie
    \item Le soglie sono associate ai singoli sensori e possono quindi variare per edificio e grandezza monitorata
    \item Al superamento di una soglia il sistema genera un allarme automatico e invia una notifica email
\end{itemize}

\paragraph*{TASKS – User Story 8:}
\begin{enumerate}
    \item \textbf{Creare l'interfaccia di configurazione soglie}
    \begin{itemize}
        \item Implementare form con campi per la soglia minima e massima del sensore
        \item Applicare le soglie a livello di singolo sensore
    \end{itemize}
    \item \textbf{Implementare la logica di validazione}
    \begin{itemize}
        \item Verificare che i valori inseriti siano coerenti
        \item Validare i range di valori accettabili
    \end{itemize}
    \item \textbf{Salvare le configurazioni}
    \begin{itemize}
        \item Memorizzare le soglie nel database associate al sensore
        \item Rimuovere le soglie e gli allarmi associati quando vengono azzerate
    \end{itemize}
    \item \textbf{Testare il sistema di soglie}
    \begin{itemize}
        \item Verificare che le soglie vengano applicate correttamente
        \item Testare la generazione di allarmi al superamento soglie
    \end{itemize}
\end{enumerate}


\subsection{User Story 9 – Confronto delle performance tra edifici}
\label{us:analisiComparative}

\textbf{Confronto performance tra edifici} \\
Come amministratore comunale, voglio confrontare le performance energetiche di diversi edifici affiancando i loro indicatori, in modo da identificare gli edifici più e meno efficienti.

\paragraph*{Criteri di accettazione:}
\begin{itemize}
    \item L'utente può selezionare più edifici da confrontare (fino a 3)
    \item Vengono mostrati affiancati gli indicatori di efficienza calcolati per ciascun edificio (energia consumata, dispersione termica stimata, COP)
    \item Il confronto si basa sulle metriche calcolate dal backend per ogni singolo edificio (efficienza termica nel periodo selezionato)
    \item Il sistema facilita l'identificazione degli edifici più e meno efficienti
\end{itemize}

\paragraph*{TASKS – User Story 9:}
\begin{enumerate}
    \item \textbf{Creare l'interfaccia di selezione edifici}
    \begin{itemize}
        \item Implementare controlli per selezionare gli edifici da confrontare (massimo 3)
        \item Mostrare le informazioni base degli edifici selezionati
    \end{itemize}
    \item \textbf{Implementare le visualizzazioni comparative}
    \begin{itemize}
        \item Creare grafici che mostrano gli indicatori di efficienza dei diversi edifici affiancati
        \item Utilizzare colori distinti per ogni edificio
    \end{itemize}
    \item \textbf{Utilizzare gli indicatori di performance}
    \begin{itemize}
        \item Recuperare le metriche di efficienza per ciascun edificio dall'endpoint dedicato
        \item Evidenziare gli edifici più e meno efficienti con indicatori visivi
    \end{itemize}
    \item \textbf{Verificare il comportamento del confronto}
    \begin{itemize}
        \item Testare il confronto con periodi temporali diversi
        \item Testare la gestione di edifici privi di dati di efficienza
    \end{itemize}
\end{enumerate}


\subsection{User Story 10 – Gestione edifici e configurazione sensori}
\label{us:aggiuntaEdificio}

\textbf{Aggiunta nuovo edificio al sistema} \\
Come operatore tecnico/amministratore, voglio aggiungere un nuovo edificio al sistema inserendo tutti i dati anagrafici e associando i sensori, in modo da iniziare a monitorare i consumi energetici.

\paragraph*{Criteri di accettazione:}
\begin{itemize}
    \item L'utente può inserire nome, indirizzo, superficie e tipologia impianto
    \item È possibile associare sensori (temperatura e consumo) specificando posizione fisica
    \item Il sistema valida i dati inseriti (campi obbligatori, tipologia esistente)
    \item I sensori trasmettono misurazioni a intervalli configurabili per ciascun sensore (valore predefinito: 90 secondi)
\end{itemize}

\paragraph*{TASKS – User Story 10:}
\begin{enumerate}
    \item \textbf{Creare il form di inserimento edificio}
    \begin{itemize}
        \item Aggiungere i campi nel form per tutti le informazioni anagrafiche
        \item Aggiungere un sistema di validazione client-side per tutti i campi obbligatori
    \end{itemize}
    \item \textbf{Implementare la validazione dei dati di input}
    \begin{itemize}
        \item Verificare la presenza dei campi obbligatori e la validità dei formati
        \item Mostrare errore in caso di dati non validi
    \end{itemize}
    \item \textbf{Creare l'interfaccia di associazione sensori}
    \begin{itemize}
        \item Permettere l'aggiunta di sensori di temperatura e di consumo
        \item Raccogliere informazioni su posizione fisica e tipologia
    \end{itemize}
    \item \textbf{Configurare la frequenza trasmissione dati}
    \begin{itemize}
        \item Impostare l'intervallo di trasmissione dei sensori (valore predefinito: 90 secondi)
        \item Verificare la corretta ricezione dei dati
    \end{itemize}
\end{enumerate}


\subsection{User Story 11 – Gestione notifiche di allarme}
\label{us:notificheAllarme}

\textbf{Visualizzazione e gestione allarmi attivi} \\
Come operatore tecnico o amministratore, voglio visualizzare tutte le notifiche di allarme attive in un pannello dedicato ordinato per timestamp, in modo da gestire prioritariamente gli allarmi più recenti o critici.

\paragraph*{Criteri di accettazione:}
\begin{itemize}
    \item Tutte le notifiche attive sono elencate in un pannello dedicato
    \item Le notifiche sono ordinate per timestamp (più recenti prime)
    \item È possibile visualizzare i dettagli di ogni allarme (edificio, valori rilevati, soglie)
    \item L'operatore può prendere in carico le notifiche e chiuderle dopo la risoluzione
    \item Il sistema registra l'operatore e il timestamp di presa in carico e di risoluzione
\end{itemize}

\paragraph*{TASKS – User Story 11:}
\begin{enumerate}
    \item \textbf{Creare il pannello notifiche allarmi}
    \begin{itemize}
        \item Implementare lista ordinata per timestamp
        \item Utilizzare indicatori visivi per priorità (colori, icone)
    \end{itemize}
    \item \textbf{Implementare la visualizzazione dettagli}
    \begin{itemize}
        \item Creare modal o pannello espandibile per i dettagli allarme
        \item Mostrare edificio coinvolto, tipo di allarme, valori rilevati e soglie
    \end{itemize}
    \item \textbf{Implementare la funzione di chiusura allarme}
    \begin{itemize}
        \item Aggiungere pulsante per prendere in carico ogni notifica (stato \texttt{acknowledged})
        \item Aggiungere pulsante per risolvere l'allarme (stato \texttt{resolved})
    \end{itemize}
    \item \textbf{Registrare le operazioni di gestione}
    \begin{itemize}
        \item Registrare operatore e timestamp di presa in carico e di risoluzione sull'allarme
        \item Gli allarmi chiusi restano consultabili tramite il filtro per stato
    \end{itemize}
    \item \textbf{Implementare filtri e ricerca}
    \begin{itemize}
        \item Permettere filtro per stato (attivi, presi in carico, risolti)
        \item Aggiungere ricerca per edificio specifico
    \end{itemize}
\end{enumerate}


\subsection{User Story 12 – Visualizzazione stato sensori}
\label{us:statoSensori}

\textbf{Monitoraggio stato di tutti i sensori} \\
Come operatore tecnico/amministratore, voglio visualizzare lo stato operativo di tutti i sensori installati negli edifici con indicazione se sono attivi o offline, in modo da identificare rapidamente eventuali malfunzionamenti.

\paragraph*{Criteri di accettazione:}
\begin{itemize}
    \item Viene mostrato un elenco completo dei sensori installati (temperatura e consumo)
    \item Per ogni sensore è indicato lo stato (attivo/offline)
    \item È visibile la data dell'ultimo aggiornamento dati
    \item I sensori offline sono evidenziati visivamente
\end{itemize}

\paragraph*{TASKS – User Story 12:}
\begin{enumerate}
    \item \textbf{Creare la visualizzazione lista sensori}
    \begin{itemize}
        \item Implementare tabella o griglia con tutti i sensori organizzati per edificio
    \end{itemize}
    \item \textbf{Implementare gli indicatori di stato}
    \begin{itemize}
        \item Utilizzare icone o colori per stato attivo/offline
        \item Evidenziare in modo prominente i sensori offline
    \end{itemize}
    \item \textbf{Mostrare informazioni dettagliate}
    \begin{itemize}
        \item Includere l'edificio di appartenenza
        \item Mostrare timestamp ultimo aggiornamento
    \end{itemize}
    \item \textbf{Implementare sistema di verifica stato}
    \begin{itemize}
        \item Controllare periodicamente la ricezione dati dai sensori
        \item Aggiornare automaticamente lo stato in base all'attività
    \end{itemize}
    \item \textbf{Aggiungere funzionalità di filtro e ricerca}
    \begin{itemize}
        \item Permettere filtro per stato (attivi, offline, tutti)
        \item Implementare ricerca per edificio
    \end{itemize}
\end{enumerate}


\subsection{User Story 13 – Monitoraggio tempo reale}
\label{us:monitoraggioTempoReale}

\textbf{Visualizzazione dati in tempo reale con aggiornamento automatico} \\
Come operatore tecnico/amministratore, voglio monitorare in tempo reale i dati di ogni edificio con aggiornamento automatico ogni 90 secondi, in modo da avere sempre i valori più recenti senza dover ricaricare manualmente la pagina.

\paragraph*{Criteri di accettazione:}
\begin{itemize}
    \item I dati si aggiornano automaticamente ogni 90 secondi (intervallo configurabile)
    \item Vengono mostrati temperatura interna, esterna e consumo istantaneo
    \item È presente un grafico delle ultime 24 ore
    \item L'invio di notifiche email automatiche per variazioni significative è implementato per gli allarmi generati dal sistema (superamento delle soglie e calo di efficienza energetica);
\end{itemize}

\paragraph*{TASKS – User Story 13:}
\begin{enumerate}
    \item \textbf{Implementare l'aggiornamento automatico}
    \begin{itemize}
        \item Configurare polling ogni 90 secondi per recuperare nuovi dati
        \item Aggiornare i valori visualizzati senza reload della pagina
    \end{itemize}
    \item \textbf{Visualizzare i parametri richiesti}
    \begin{itemize}
        \item Mostrare temperatura interna, esterna e consumo con indicatori chiari
        \item Utilizzare visualizzazioni appropriate (gauge, numeri grandi, ecc.)
    \end{itemize}
    \item \textbf{Implementare il grafico delle 24 ore}
    \begin{itemize}
        \item Creare grafico a linee con andamento temporale
        \item Aggiornare il grafico con i nuovi dati
    \end{itemize}
    \item \textbf{Segnalare i malfunzionamenti dei sensori}
    \begin{itemize}
        \item Evidenziare nell'interfaccia i sensori offline o in errore
        \item Indicare la data dell'ultima rilevazione per ogni sensore
    \end{itemize}
\end{enumerate}


\subsection{User Story 14 – Gestione edifici (rimozione)}
\label{us:rimozioneEdificio}

\textbf{Rimozione edificio dismesso} \\
Come operatore tecnico/amministratore, voglio poter rimuovere dal sistema gli edifici che sono stati dismessi o non sono più monitorati, in modo da mantenere il database aggiornato e contenere solo informazioni attuali.

\paragraph*{Criteri di accettazione:}
\begin{itemize}
    \item L'utente può selezionare un edificio da rimuovere
    \item Il sistema richiede conferma prima della rimozione definitiva
    \item Vengono rimossi a cascata anche i sensori, gli allarmi e le rilevazioni associati
    \item L'operazione è consentita agli utenti con ruolo amministratore o operatore
\end{itemize}

\paragraph*{TASKS – User Story 14:}
\begin{enumerate}
    \item \textbf{Creare l'interfaccia di rimozione}
    \begin{itemize}
        \item Aggiungere pulsante ``Rimuovi edificio'' nella pagina di dettaglio
        \item Implementare lista edifici con opzione di rimozione
    \end{itemize}
    \item \textbf{Implementare la conferma rimozione}
    \begin{itemize}
        \item Creare modal di conferma con dettagli edificio
        \item Richiedere conferma esplicita dall'operatore
    \end{itemize}
    \item \textbf{Gestire la cancellazione a cascata}
    \begin{itemize}
        \item Eliminare le rilevazioni associate all'edificio
        \item Eliminare tutti gli allarmi e i sensori associati
    \end{itemize}
    \item \textbf{Rimuovere edificio e sensori associati}
    \begin{itemize}
        \item Eliminare l'edificio dal database attivo
        \item Gestire correttamente le relazioni nel database
    \end{itemize}
\end{enumerate}


\subsection{User Story 15 – Gestione edifici (modifica)}
\label{us:modificaEdificio}

\textbf{Modifica edificio} \\
Come operatore tecnico/amministratore, voglio poter modificare nel sistema un edificio che ha subito degli aggiornamenti, in modo da mantenere il database aggiornato e contenere solo informazioni attuali.

\paragraph*{Criteri di accettazione:}
\begin{itemize}
    \item L'utente può modificare l'edificio selezionato
    \item Il sistema valida i dati prima del salvataggio
    \item Vengono mantenuti tutti i sensori associati
    \item Lo storico dei dati viene preservato
    \item Vengono registrati sia l'autore della modifica sia il momento della modifica 
\end{itemize}

\paragraph*{TASKS – User Story 15:}
\begin{enumerate}
    \item \textbf{Creare l'interfaccia di modifica}
    \begin{itemize}
        \item Aggiungere pulsante ``Modifica edificio'' nella pagina di dettaglio
    \end{itemize}
    \item \textbf{Implementare la modifica}
    \begin{itemize}
        \item Creare modal da compilare con i dettagli dell'edificio
        \item Richiedere conferma esplicita dall'operatore
    \end{itemize}
    \item \textbf{Gestire l'aggiornamento dei dati}
    \begin{itemize}
        \item Aggiornare l'edificio dal database attivo
        \item Mantenere tutti i sensori associati
        \item Gestire correttamente le relazioni nel database
    \end{itemize}
    \item \textbf{Tracciare la modifica}
    \begin{itemize}
        \item Salvare l'operatore che ha effettuato la modifica
        \item Aggiornare il timestamp di quando è stata effettuata la modifica
    \end{itemize}
\end{enumerate}

\subsection{User Story 16 – Accettazione invito e primo accesso}
\label{us:accettazioneInvito}

\textbf{Completamento della registrazione tramite invito} \\
Come utente invitato alla piattaforma, voglio completare la registrazione impostando la mia password tramite il link ricevuto, in modo da poter accedere al sistema con il ruolo che mi è stato assegnato.

\paragraph*{Criteri di accettazione:}
\begin{itemize}
    \item Il link di invito è univoco e ha scadenza 7 giorni dopo l'invio
    \item L'utente può scegliere tra la creazione dell'account con password o tramite l'uso del Single Sign-On di Google
    \item Il sistema mostra un messaggio di errore quando il token è scaduto, revocato o già utilizzato
    \item Dopo la registrazione, l'utente viene autenticato e reindirizzato alla dashboard principale
    \item Nell'invito viene specificato anche il ruolo assegnato
\end{itemize}

\paragraph*{TASKS – User Story 16:}
\begin{enumerate}
    \item \textbf{Creare la pagina di accettazione invito}
    \begin{itemize}
        \item Implementare la pagina di registrazione accessibile solo tramite link
        \item Mostrare email e ruolo associati all'invito
    \end{itemize}
    \item \textbf{Validare il token di invito}
    \begin{itemize}
        \item Verificare validità, scadenza e stato dell'invito
        \item Gestire gli inviti scaduti, revocati o già accettati
    \end{itemize}
    \item \textbf{Gestire le due modalità di registrazione}
    \begin{itemize}
        \item Implementare il form di impostazione password e nome
        \item Implementare l'accettazione dell'invito tramite Google
    \end{itemize}
    \item \textbf{Completare la creazione dell'account}
    \begin{itemize}
        \item Creare l'account con il ruolo dell'invito
        \item Autenticare l'utente e reindirizzarlo alla dashboard
    \end{itemize}
\end{enumerate}


\subsection{User Story 17 – Reset della password}
\label{us:resetPassword}

\textbf{Recupero dell'accesso tramite password dimenticata} \\
Come utente, voglio recuperare l'accesso al mio account tramite la procedura "Password dimenticata", in modo da reimpostare la mia password in autonomia.

\paragraph*{Criteri di accettazione:}
\begin{itemize}
    \item L'utente può richiedere il reset inserendo la propria email
    \item Il sistema invia un link di reset con token monouso e scadenza (1 ora)
    \item La risposta alla richiesta è generica per evitare l'enumerazione delle email
    \item Il token può essere utilizzato una sola volta
    \item Dopo il reset l'utente può accedere con la nuova password
\end{itemize}

\paragraph*{TASKS – User Story 17:}
\begin{enumerate}
    \item \textbf{Implementare la richiesta di reset}
    \begin{itemize}
        \item Creare la pagina di inserimento email
        \item Generare un token sicuro, memorizzato come hash, con scadenza
    \end{itemize}
    \item \textbf{Implementare l'invio email}
    \begin{itemize}
        \item Configurare il template email con il link di reset
        \item Inviare l'email solo se l'account esiste
    \end{itemize}
    \item \textbf{Implementare la validazione e il reset}
    \begin{itemize}
        \item Validare il token prima del reset
        \item Aggiornare la password e invalidare il token dopo l'uso
    \end{itemize}
\end{enumerate}


\subsection{User Story 18 – Soglie di efficienza energetica (COP)}
\label{us:efficienzaCop}

\textbf{Configurazione della soglia minima di efficienza} \\
Come operatore tecnico o amministratore, voglio configurare la soglia minima di efficienza (COP) per un edificio, in modo da ricevere allarmi automatici quando l'efficienza energetica scende sotto soglia.

\paragraph*{Criteri di accettazione:}
\begin{itemize}
    \item L'utente può attivare o disattivare la soglia di efficienza per ciascun edificio
    \item È possibile definire il valore minimo di COP accettabile
    \item Il sistema valuta periodicamente l'efficienza dell'edificio su una finestra di 24 ore
    \item Se l'efficienza scende sotto soglia, il sistema genera un allarme automatico e invia una notifica email
    \item La soglia non è disponibile per gli edifici a teleriscaldamento
    \item Gli allarmi di efficienza vengono risolti automaticamente al rientro sopra soglia
\end{itemize}

\paragraph*{TASKS – User Story 18:}
\begin{enumerate}
    \item \textbf{Creare l'interfaccia di configurazione efficienza}
    \begin{itemize}
        \item Implementare toggle di attivazione e campo per la soglia minima di COP
        \item Escludere la configurazione per gli edifici a teleriscaldamento
    \end{itemize}
    \item \textbf{Implementare la valutazione periodica}
    \begin{itemize}
        \item Calcolare l'efficienza media su una finestra di 24 ore
        \item Confrontare il valore calcolato con la soglia configurata
    \end{itemize}
    \item \textbf{Gestire gli allarmi di efficienza}
    \begin{itemize}
        \item Generare l'allarme al superamento di una soglia
        \item Chiusura automatica degli allarmi al rientro dentro la soglia
    \end{itemize}
    \item \textbf{Testare il sistema di soglie di efficienza}
    \begin{itemize}
        \item Verificare la generazione e la risoluzione degli allarmi
        \item Testare gli edifici privi di dati sufficienti
    \end{itemize}
\end{enumerate}


\subsection{User Story 19 – Visualizzazione edifici su mappa}
\label{us:mappaEdifici}

\textbf{Visualizzazione degli edifici su mappa interattiva} \\
Come operatore tecnico o amministratore, voglio visualizzare gli edifici comunali monitorati su una mappa interattiva, in modo da identificare rapidamente la dislocazione sul territorio e lo stato delle strutture.

\paragraph*{Criteri di accettazione:}
\begin{itemize}
    \item La mappa mostra un indicatore per ciascun edificio monitorato sul territorio comunale.
    \item Ogni indicatore è differenziato cromaticamente in base allo stato operativo dello stabile (verde per attivo, giallo per inattivo, rosso per dismesso).
    \item Cliccando su un indicatore viene mostrato un riepilogo sintetico (nome, indirizzo, stato) con pulsante di collegamento diretto alla pagina di dettaglio dell'edificio.
    \item Le coordinate geografiche vengono determinate automaticamente dall'indirizzo tramite un servizio di geocoding, con anche la possibilità di cambiare manuale in fase di gestione edificio.
\end{itemize}

\paragraph*{TASKS – User Story 19:}
\begin{enumerate}
    \item \textbf{Creare la vista cartografica}
    \begin{itemize}
        \item Integrare il componente mappa interattiva nel layout applicativo.
        \item Posizionare i marker in base alle coordinate memorizzate nel database.
    \end{itemize}
    \item \textbf{Implementare la legenda e gli indicatori di stato}
    \begin{itemize}
        \item Associare la codifica cromatica ai marker conformi alla legenda (attivo, inattivo, dismesso).
        \item Configurare i popup informativi al passaggio del mouse o al click.
    \end{itemize}
    \item \textbf{Gestire geocoding e navigazione}
    \begin{itemize}
        \item Integrare il geocoding automatico dell'indirizzo e la gestione delle eccezioni per indirizzi non localizzabili.
        \item Aggiungere il link di navigazione diretta dal marker alla dashboard di dettaglio edificio.
    \end{itemize}
\end{enumerate}

\subsection{User Story 20 – Disabilitazione account utente}
\label{us:disabilitazioneAccount}

\textbf{Disabilitazione e riabilitazione temporanea degli account} \\
Come amministratore comunale, voglio poter disabilitare o riabilitare temporaneamente l'accesso di un utente senza cancellarne l'account, in modo da revocare tempestivamente i permessi preservando così lo storico delle attività svolte.

\paragraph*{Criteri di accettazione:}
\begin{itemize}
    \item L'amministratore può modificare lo stato di un account (attivo/disabilitato) direttamente dalla tabella di gestione utenti.
    \item Un utente con account disabilitato non può effettuare il login né tramite credenziali locali né tramite Google (il sistema mostra un chiaro avviso di account sospeso).
    \item La disabilitazione invalida le eventuali sessioni attive dell'utente coinvolto.
    \item L'amministratore non può in alcun caso disabilitare, revocare o eliminare il proprio account personale.
    \item Gli account disabilitati sono chiaramente riconoscibili visivamente nella lista utenti.
    \item Lo storico delle modifiche pregresse resta associato all'utente anche in stato disabilitato.
\end{itemize}

\paragraph*{TASKS – User Story 20:}
\begin{enumerate}
    \item \textbf{Adeguamento del modello dati utente}
    \begin{itemize}
        \item Introdurre il campo di stato operativo nel modello dell'utente.
        \item Implementare gestione utenti nelle API per consentire la commutazione dello stato.
    \end{itemize}
    \item \textbf{Interfaccia di amministrazione utenti}
    \begin{itemize}
        \item Inserire nella tabella utenti il pulsante per disabilitare o riabilitare l'utenza.
        \item Evidenziare graficamente gli utenti disabilitati rispetto a quelli attivi.
    \end{itemize}
    \item \textbf{Vincoli di sicurezza e controllo accessi}
    \begin{itemize}
        \item Verificare lo stato dell'account nei middleware di autenticazione (locale e Google) bloccando gli utenti disabilitati.
        \item Impedire via codice e interfaccia l'auto-disabilitazione dell'amministratore collegato.
    \end{itemize}
    \item \textbf{Revoca immediata delle sessioni}
    \begin{itemize}
        \item Invalidare i token JWT/sessioni attive al momento della disabilitazione.
    \end{itemize}
\end{enumerate}